\documentclass{beamer}
\usepackage{lmodern}

\usepackage{xeCJK}

\usepackage[safe]{silence}
\WarningFilter*{latexfont}{Font}

\usetheme{CambridgeUS} % try Madrid, Pittsburgh
\usecolortheme{beaver}
\usefonttheme[onlymath]{serif} % try "professionalfonts"

\setbeamertemplate{itemize items}[default]
\setbeamertemplate{enumerate items}[default]

\usepackage{amsmath, amsfonts, latexsym, mathtools}
% \usepackage{kbordermatrix}
% \usepackage{ntheorem}

\usepackage{centernot}

\DeclareMathOperator*{\argmin}{arg\,min}
\DeclareMathOperator*{\argmax}{arg\,max}

% colors
\newcommand{\red}[1]{\textcolor{red}{#1}}
\newcommand{\green}[1]{\textcolor{green}{#1}}
\newcommand{\blue}[1]{\textcolor{blue}{#1}}
\newcommand{\purple}[1]{\textcolor{purple}{#1}}

% color box
\newcommand{\rbox}[1]{\red{\boxed{#1}}}
\newcommand{\gbox}[1]{\green{\boxed{#1}}}
\newcommand{\bbox}[1]{\blue{\boxed{#1}}}
\newcommand{\pbox}[1]{\purple{\boxed{#1}}}

\usepackage{tikz}
% see http://tex.stackexchange.com/a/7045/23098
\newcommand*\circled[1]{\tikz[baseline=(char.base)]{
            \node[shape=circle,draw,inner sep=2pt] (char) {#1};}}
\usepackage{tikz-qtree}
\usetikzlibrary{backgrounds, fit}

\usepackage{pifont}
\usepackage{wasysym}

\usepackage[normalem]{ulem}
\newcommand{\middlewave}[1]{\raisebox{0.5em}{\uwave{\hspace{#1}}}}

\usepackage{graphicx, subcaption}

\usepackage{algorithm}
\usepackage[noend]{algpseudocode}

\newcommand{\pno}[1]{\textcolor{blue}{\scriptsize [Problem: #1]}}
\newcommand{\set}[1]{\{#1\}}
\DeclareMathOperator{\lcm}{lcm}

\newcommand{\cmark}{\green{\ding{51}}}
\newcommand{\xmark}{\red{\ding{55}}}
%%%%%%%%%%%%%%%%%%%%%%%%%%%%%%%%%%%%%%%%%%%%%%%%%%%%%%%%%%%%%%
% for fig without caption: #1: width/size; #2: fig file
\newcommand{\fignocaption}[2]{
  \begin{figure}[htp]
    \centering
      \includegraphics[#1]{#2}
  \end{figure}
}

% for fig with caption: #1: width/size; #2: fig file; #3: fig caption
\newcommand{\fig}[3]{
  \begin{figure}[htp]
    \centering
      \includegraphics[#1]{#2}
      \caption[labelInTOC]{#3}
  \end{figure}
}

\newcommand{\titletext}{Algebraic Coding Theory (Revisited)}
%%%%%%%%%%%%%%%%%%%%
\title[\titletext]{\titletext}
\subtitle{--- From the perspective of linear algebra}

\author[Hengfeng Wei]{Hengfeng Wei}
\titlegraphic{\includegraphics[height = 1.5cm]{figs/qrcode-coding-theory-20170420.png}}
\institute{hfwei@nju.edu.cn}
\date{April 17 $\sim$ April 20, 2017}

\AtBeginSection[]{
  \begin{frame}[noframenumbering, plain]
    \frametitle{\titletext}
    \tableofcontents[currentsection, sectionstyle=show/shaded, subsectionstyle=show/show/hide]
  \end{frame}
}
%%%%%%%%%%
\begin{document}
\maketitle

\section{Block Codes}

%%%%%%%%%%%%%%%%%%%%
\begin{frame}{Block coding}
  \centerline{flow chart here}

  \centerline{Why is cryptography?}
\end{frame}
%%%%%%%%%%%%%%%%%%%%
\begin{frame}{Important code parameters}
  \[
	k
  \]

  \[
	n > k
  \]

  \[
	r = n - k
  \]

  \[
	|C| \le 2^n
  \]

  \[
	0 < R \triangleq \frac{k}{n} < 1
  \]
\end{frame}
%%%%%%%%%%%%%%%%%%%%
\begin{frame}{Hamming distance}
  \[
	w(c) = \# 1's \text{ in } c
  \]

  \[
	d(c_1, c_2) = w(c_1 + c_2)
  \]

  \[
	\begin{aligned}
	  d(C) &= \min \set{d(c_1, c_2) \mid c_1 \neq c_2, c_1, c_2 \in C} \\
	  &= \min \set{w(c_1 + c_2) \mid c_1 \neq c_2, c_1, c_2 \in C} \\
	  &\red{\neq} \min \set{w(c) \mid c \neq 0, c \in C}
	\end{aligned}
  \]
\end{frame}
%%%%%%%%%%%%%%%%%%%%
\begin{frame}{Detecting and correcting errors}
  \[
	d(C) \ge 2t + 1 \implies 2t\text{-detecting}
  \]

  \[
	d(C) \ge 2t + 1 \implies t\text{-correcting}
  \]
\end{frame}
%%%%%%%%%%%%%%%%%%%%
\begin{frame}{Sphere-packing bound}
  \begin{theorem}[Sphere-packing bound]
	A $t$-error-correcting binary code of length $n$ must satisfy

	\[
	  |C| \sum_{i=0}^{t} \binom{n}{i} \le 2^n
	\]
  \end{theorem}

  \[
	t = 1 \implies |C| \le \frac{2^n}{n + 1}
  \]

  \begin{definition}[Perfect code]
	\[
	  |C| \sum_{i=0}^{t} \binom{n}{i} = 2^n
	\]
  \end{definition}
\end{frame}
%%%%%%%%%%%%%%%%%%%%

\section{Linear Codes}

%%%%%%%%%%%%%%%%%%%%
\begin{frame}{Definition of linear code}
  \begin{definition}[Linear code]
	A linear code $C$ of length $n$ is a linear subspace of the vector space $\mathbb{F}_2^n$.
  \end{definition}

  \[
	c_1 \in C, c_2 \in C \implies c_1 + c_2 \in C
  \]

  \[
	\begin{aligned}
	  d(C) &= \min \set{w(c_1 + c_2) \mid c_1 \neq c_2, c_1, c_2 \in C} \\
	  &{\blue{=}} \min \set{w(c) \mid c \neq 0, c \in C}
	\end{aligned}
  \]
\end{frame}
%%%%%%%%%%%%%%%%%%%%
\begin{frame}{Definition of linear code}
  \begin{exampleblock}{Problem TJ--8.18}
	Let $C$ be a linear code.

	Show that either the $i$-th coordinates in the codewords of $C$ are all zeros
	or exactly half of them are zeros.
  \end{exampleblock}
\end{frame}
%%%%%%%%%%%%%%%%%%%%
\begin{frame}{Definition of linear code}
  \begin{exampleblock}{Problem TJ--8.19}
	Let $C$ be a linear code.

	Show that either every codeword has even weight
	or exactly half of them have even weight.
  \end{exampleblock}

  \[
	\text{Parity: } w(c_1) + w(c_2) \emph{ vs. } w(c_1 + c_2)
  \]
\end{frame}
%%%%%%%%%%%%%%%%%%%%
\begin{frame}{Definition of linear code}
  \begin{definition}[Linear code]
	An $(n,k)$ linear code $C$ of length $n$ and rank $k$ is a linear subspace with dimension $k$ of the vector space $\mathbb{F}_2^n$.
  \end{definition}

  \[
	\text{Basis: } c_1, c_2, \dots, c_k
  \]

  \[
	c_i = \alpha_1 c_1 + \alpha_2 c_2 + \cdots + \alpha_k c_k
  \]

  \[
	|C| = 2^k
  \]
\end{frame}
%%%%%%%%%%%%%%%%%%%%
\begin{frame}{Generator matrix}
  \begin{definition}[Generator matrix]
	A matrix $G_{n \times k}$ is a generator matrix for an $(n,k)$ linear code $C$ if 
	
	\[
	  C = \text{Col}(G)
	\]
  \end{definition}

  \[
	\text{rank}(G) = k
  \]

  \[
	G_{(n \times k)} \cdot d_{k \times 1} = c_{n \times 1} \in C
  \]

  \[
	G(c_1 + c_2) = G(c_1) + G(c_2)
  \]
\end{frame}
%%%%%%%%%%%%%%%%%%%%
\begin{frame}{Generator matrix for Hamming code $(7,4)$}
  \begin{columns}
	\column{0.45\textwidth}
	  \fignocaption{width = 0.60\textwidth}{figs/hamming74-venn.png}
	\column{0.45\textwidth}
	  \[
		G = \begin{bmatrix}
		  1 & 0 & 0 & 0 \\
		  0 & 1 & 0 & 0 \\
		  0 & 0 & 1 & 0 \\
		  0 & 0 & 0 & 1 \\
		  1 & 1 & 0 & 1 \\
		  0 & 1 & 1 & 1 \\
		  1 & 0 & 1 & 1 \\
		\end{bmatrix}
	  \]

  \end{columns}

  \[
	G \cdot \begin{pmatrix}
	  1 \\ 0 \\ 1 \\ 0
	\end{pmatrix}
	= 
  \]
\end{frame}
%%%%%%%%%%%%%%%%%%%%
\begin{frame}{Standard generator matrix}
  \begin{exampleblock}{Problem TJ--8.7}
	Generator matrices are NOT unique.
  \end{exampleblock}

  \begin{definition}[Standard generator matrix]
	A generator matrix $G_{n \times k}$ is standard if

	\[
	  G_{n \times k} = \begin{bmatrix}
		I_k \\ A_{(n-k) \times k}
	  \end{bmatrix}
	\]
  \end{definition}
\end{frame}
%%%%%%%%%%%%%%%%%%%%
\begin{frame}{From generator matrix to parity-check matrix}
  \[
	G \cdot \begin{pmatrix}
      d_1 \\ d_2 \\ d_3 \\ d_4
	\end{pmatrix}
	= \begin{bmatrix}
	  1 & 0 & 0 & 0 \\
	  0 & 1 & 0 & 0 \\
	  0 & 0 & 1 & 0 \\
	  0 & 0 & 0 & 1 \\
	  1 & 1 & 0 & 1 \\
	  0 & 1 & 1 & 1 \\
	  1 & 0 & 1 & 1 \\
	\end{bmatrix} 
	\cdot \begin{pmatrix}
      d_1 \\ d_2 \\ d_3 \\ d_4
	\end{pmatrix}
	= \begin{pmatrix}
      d_1 \\ d_2 \\ d_3 \\ d_4 \\
	  p_1 = d_1 + d_2  \qquad + d_4 \\
	  p_2 = \qquad d_2 + d_3 + d_4 \\
	  p_3 = d_1 \qquad + d_3 + d_4 \\
	\end{pmatrix}
  \]
\end{frame}
%%%%%%%%%%%%%%%%%%%%
\begin{frame}{From generator matrix to parity-check matrix}
  \begin{gather*}
	d_1 + d_2  \qquad + d_4 + p_1 \qquad\qquad = 0 \\
	\qquad d_2 + d_3 + d_4 \qquad + p_2 \qquad = 0 \\
	d_1 \qquad + d_3 + d_4 \qquad \qquad + p_3 = 0 \\
  \end{gather*}

  \[
	\begin{bmatrix}
	  1 & 1 & 0 & 1 & 1 & 0 & 0 \\
	  0 & 1 & 1 & 1 & 0 & 1 & 0 \\
	  1 & 0 & 1 & 1 & 0 & 0 & 1 \\
	\end{bmatrix} \begin{pmatrix}
      d_1 \\ d_2 \\ d_3 \\ d_4 \\ p_1 \\ p_2 \\ p_3
	\end{pmatrix}
	= 0
  \]
\end{frame}
%%%%%%%%%%%%%%%%%%%%
\begin{frame}{Parity-check matrix}
  \begin{definition}[Parity-check matrix]
	A matrix $H_{(n - k) \times n}$ is a parity-check matrix for an $(n,k)$ linear code $C$ if

	\[
	  C = \text{Nul}(H)
	\]
  \end{definition}

  \[
	H_{(n - k) \times n} \cdot c_{n \times 1} = 0_{(n - k) \times 1}
  \]

  \[
	\text{rank}(H) = n - k
  \]
\end{frame}
%%%%%%%%%%%%%%%%%%%%
\begin{frame}{Standard parity-check matrix}
  \begin{exampleblock}{Problem TJ--8.11}
	Parity-check matrices are NOT unique.
  \end{exampleblock}

  \centerline{Elementary row operations.}

  \vspace{0.60cm}
  \begin{definition}[Standard parity-check matrix]
	A parity-check matrix $H_{(n-k) \times n}$ is standard if

	\[
	  H_{(n-k) \times n} = \begin{bmatrix}
		A_{(n-k) \times k} \mid I_{n-k}
	  \end{bmatrix}
	\]
  \end{definition}
\end{frame}
%%%%%%%%%%%%%%%%%%%%
\begin{frame}{Generator matrix and Parity-check matrix}
  \[
	H_{(n-k) \times n} \cdot G_{n \times k} \cdot d_{k \times 1} = 0_{(n-k) \times 1}
  \]

  \begin{align*}
	H_{(n-k) \times n} &\cdot G_{n \times k} \\
	&= \begin{bmatrix}
	  A_{(n-k) \times k} \mid I_{n-k} 
	\end{bmatrix}
	\cdot
	\begin{bmatrix}
	  I_k \\ A_{(n-k) \times k}
	\end{bmatrix} \\
	&= A_{(n-k) \times k} \cdot I_k + I_{n-k} \cdot A_{(n-k) \times k} \\
	&= A_{(n-k) \times k} + A_{(n-k) \times k} \\
	&= 0_{(n-k) \times k}
  \end{align*}
\end{frame}
%%%%%%%%%%%%%%%%%%%%
\begin{frame}{Syndrome decoding}
  \[
	r = c + e_i
  \]

  \[
	r = c + (e_{i} + e_{j} + \cdots)
  \]

  \begin{definition}[Syndrome]
	\begin{align*}
	  S(r) &= H r \\
	  	&= H(c + (e_{i} + e_{j} + \cdots)) \\
		&= H (e_{i} + e_{j} + \cdots) \\
		&= H e_{i} + H e_{j} + \cdots \\
		&= S(e_i) + S{e_{j}} + \cdots
	\end{align*}
  \end{definition}
\end{frame}
%%%%%%%%%%%%%%%%%%%%
\begin{frame}{Syndrome decoding}
  \begin{exampleblock}{Problem TJ--8.13}
	\[
	  H = \begin{bmatrix}
		0 & 1 & 1 & 1 & 1 \\
		0 & 0 & 0 & 1 & 1 \\
		1 & 0 & 1 & 0 & 1
	  \end{bmatrix}
	\]

	\centerline{(d) Errors in the third and fourth bits}
  \end{exampleblock}
\end{frame}
%%%%%%%%%%%%%%%%%%%%
\begin{frame}{Extracting $d(C)$ from $H$}
  \begin{theorem}[]
	If $H$ is the parity-check matrix for a linear code $C$, then
	$d(C)$ equals the minimum number of columns of $H$ that are linearly dependent.
  \end{theorem}

  \[
	\begin{bmatrix}
	  1 & 0 & 1 & 0 \\
	  0 & 1 & 1 & 0 \\
	  0 & 0 & 0 & 1
	\end{bmatrix}
  \]
\end{frame}
%%%%%%%%%%%%%%%%%%%%
\begin{frame}{Error-detecting and error-correcting capabilities}
  \begin{align*}
	t = 1 & \implies d(C) = 3 \\
	  &\iff \forall \set{c_i}, \forall \set{c_i, c_j} \text{ linearly independent} \\
	  &\iff \text{no zero column, no identical columns}
  \end{align*}
\end{frame}
%%%%%%%%%%%%%%%%%%%%
\begin{frame}{Error-detecting and error-correcting capabilities}
  \begin{exampleblock}{Problem TJ--8.21; TJ--8.23}
	If we are to use an error-correcting linear code to transmit the 128 ASCII characters,
	what size matrix must be used? What if we require only error detection?
  \end{exampleblock}

  \[
	t = 1
  \]

  \[
	2^r - 1 \ge 7 + r \implies r \ge 4
  \]
\end{frame}
%%%%%%%%%%%%%%%%%%%%
\begin{frame}{Generalized Hamming codes}
  \[
	C_{Ham} = (7,4)
  \]

  \[
	H_{3 \times 7}: r = 3, n = 2^3 - 1 = 7
  \]

  \[
	\begin{bmatrix}
	  1 & 1 & 0 & 1 & 1 & 0 & 0 \\
	  0 & 1 & 1 & 1 & 0 & 1 & 0 \\
	  1 & 0 & 1 & 1 & 0 & 0 & 1 \\
	\end{bmatrix} \begin{pmatrix}
      d_1 \\ d_2 \\ d_3 \\ d_4 \\ p_1 \\ p_2 \\ p_3
	\end{pmatrix}
	= 0
  \]
\end{frame}
%%%%%%%%%%%%%%%%%%%%
\begin{frame}{Generalized Hamming codes}
  \begin{definition}[Generalized Hamming codes]
	\[
	  H_r \triangleq H_{r \times (2^r-1)}: \text{all nonzero binary vectors of length } r
	\]

	\[
	  C_{Ham}^{r} = (n = 2^r-1, k = 2^r-1-r)
	\]
  \end{definition}

  \[
	|C_{Ham}^{r}| = \frac{2^n}{n+1}
  \]

  \[
	R(C_{Ham}^{r}) = \frac{k}{n}
  \]
\end{frame}
%%%%%%%%%%%%%%%%%%%%
\tx{}
%%%%%%%%%%%%%%%%%%%%


\end{document}
%%%%%%%%%%
